\section{Finite computational observations}\label{app:finite-observations}

This appendix describes stored outputs of finite computations for the permanent and polynomial inequalities. It separates exact arithmetic from numerical optimization. A successful optimization run does not certify a global minimum. A small numerical constraint residual does not certify exact feasibility.

Write \(c_n=n!/n^n\), \(g_n=((n-1)/n)^{n-1}\), and

\[
\mathcal R(A)=\frac{\per(A)^2}{c_n\per(AA^{\T})}.
\]

For a balanced degree-\(n\) polynomial, write \(b=[x_1\cdots x_n]p\), \(N=N(p)\), and let \(H_j\) be its capacity-balanced singleton-deletion norm. The three polynomial ratios used below are

\[
D_{\rm av}=\frac{\sum_jH_j}{ng_nN},\qquad
D_{\max}=\frac{\max_jH_j}{g_nN},\qquad
S=\frac{b^2}{c_nN}.
\]

The proposed lower bound for each ratio is one. A raw deletion ratio uses \(N([x_j]p)\) in place of \(H_j\). All decimal values below are approximate numerical observations, except integers and fractions explicitly identified as exact.

\subsection{Exact coefficient tests outside the complete circulant orders}\label{app:finite-1}

Let \(C(w)_{ij}=w_{j-i\bmod n}\), and put \(F_n(w)=n^n\per(C(w))^2-n!\per(C(w)C(w)^{\T})\). For a sorting permutation \(\pi\), use \(w_{\pi(k)}=x_0+\cdots+x_k\). Each sorting region is called a chamber. The following selected chambers were checked by exact integer polynomial arithmetic.

{\def\LTcaptype{none} % do not increment counter
\begin{longtable}[]{@{}
  >{\raggedleft\arraybackslash}p{(\linewidth - 10\tabcolsep) * \real{0.1667}}
  >{\raggedleft\arraybackslash}p{(\linewidth - 10\tabcolsep) * \real{0.1667}}
  >{\raggedleft\arraybackslash}p{(\linewidth - 10\tabcolsep) * \real{0.1667}}
  >{\raggedleft\arraybackslash}p{(\linewidth - 10\tabcolsep) * \real{0.1667}}
  >{\raggedleft\arraybackslash}p{(\linewidth - 10\tabcolsep) * \real{0.1667}}
  >{\raggedleft\arraybackslash}p{(\linewidth - 10\tabcolsep) * \real{0.1667}}@{}}
\toprule\noalign{}
\begin{minipage}[b]{\linewidth}\raggedleft
Order
\end{minipage} & \begin{minipage}[b]{\linewidth}\raggedleft
Checked affine chamber representatives
\end{minipage} & \begin{minipage}[b]{\linewidth}\raggedleft
All affine representatives
\end{minipage} & \begin{minipage}[b]{\linewidth}\raggedleft
Positive coefficients in each checked chamber
\end{minipage} & \begin{minipage}[b]{\linewidth}\raggedleft
Zero coefficients
\end{minipage} & \begin{minipage}[b]{\linewidth}\raggedleft
Least positive coefficient
\end{minipage} \\
\midrule\noalign{}
\endhead
\bottomrule\noalign{}
\endlastfoot
8 & 16 & 1,260 & 245,149 & 8 & 16,736,896 \\
9 & 5 & 6,720 & 1,562,266 & 9 & 387,057,609 \\
\end{longtable}
}

Every checked coefficient was nonnegative. These are exact statements about the listed chambers. Neither full order was checked.

At order twelve, eight restrictions used the baseline plus two nested gap directions. Each restricted polynomial had 322 nonzero terms and no negative coefficient. The smallest coefficient over these restrictions was 8,915,621,446,656. A three-variable restriction does not certify its full chamber.

A noncyclic Latin-square array of order six was also checked in 50 selected sorting orders. Its rows, using indices zero through five, were

\[
\begin{pmatrix}
0&1&2&3&4&5\\3&4&1&2&5&0\\1&0&5&4&3&2\\
2&3&0&5&1&4\\4&5&3&0&2&1\\5&2&4&1&0&3
\end{pmatrix}.
\]

No checked coefficient was negative. The smallest positive coefficient was 42,624. This finite test does not establish an all-Latin-square theorem.

Additional exact tests used arbitrary lists of permutation matrices. Thirty lists of four permutations at order five, ten lists of six at order six, and thirty lists of four at order eight passed their recorded sorting checks. At order three, all sorting orders of the list of all six permutation matrices passed. The sampled lists do not prove the general coefficient assertion. The exact order-four negative coefficient in the main text refutes that assertion.

\subsection{Wider finite checks for the three-shift family}\label{app:finite-2}

For \(C=aP_n^{-1}+bI+cP_n\), exact transfer-polynomial expansion checked all six weight orders for every \(5\le n\le100\). Direct permanent expansion checked the six chambers at orders three and four. It also checked both transfer identities in orders one through seven.

A distinct implementation used characteristic-polynomial recurrences and Newton identities directly in gap coordinates. It checked 201 reflection-reduced chambers in orders three through sixty-nine, with 693,249 coefficient positions. No coefficient was negative. At orders four through sixty-nine, every checked coefficient was positive. The overlapping minima agreed with the transfer expansion.

A separate integer recurrence checked all seven nonzero binary triples in orders three through one hundred. There were 686 values. Every deficit at orders four through one hundred was positive. These checks establish finite coefficient and strictness statements for this three-shift family. They do not cover arbitrary circulant weights. The main proof uses only the finite range required by its analytic bound. Appendix~\ref{app:additional-results} records the separate exact certificate through order ninety-three that supplies both alternative transfer arguments.

\subsection{Matrix ratio searches}\label{app:finite-3}

Constrained local minimization tested \(\mathcal R(A)\) from 30 initial doubly stochastic matrices in each of orders five, six, and seven. All 90 runs reported convergence. Each order had a smallest recorded ratio of one, within numerical precision. Permanents and gradients used a positive-term subset dynamic program. Finite differences checked the gradient, and direct permutation enumeration checked one permanent.

A separate diagnostic tested 1,000 admissible centered rank-two pairs at each of eight orders:
\[
3,4,5,6,8,10,15,20.
\]
It used elementary-symmetric recurrences in extended floating precision. No violation of the strong rank-two target or the proposed duplication-to-Gram bound was reported. The input data, output values, and checking program are not available. These 8,000 reported cases cannot be reproduced from the retained data.

Sparse regular matrices were sampled as sums of disjoint permutation matrices. A positive-term subset dynamic program evaluated the permanents in extended floating precision. The stored aggregate observations are
\[
\begin{array}{c|c|c|c}
\text{Order}&\text{Degree}&\text{Samples}&\text{Approximate range of }\mathcal R(A)\\\hline
16&3&30&16.0\text{--}24.3\\
20&3&20&37.0\text{--}54.5\\
20&4&20&23.4\text{--}26.7\\
20&7&20&5.38\text{--}5.43
\end{array}.
\]
These tests found no counterexample. They were not exhaustive. Individual sampled matrices and detailed run outputs are unavailable. The ratios are numerical observations, not exact bounds. The data directory includes the numerical sampling program \texttt{sparse\_regular\_sampler.cpp}.

\subsection{Circulant ratio and Fourier-phase searches}\label{app:finite-4}

Full-simplex searches used analytic gradients and sequential least-squares programming, with \(w\ge0\), \(\sum w_i=1\), and at most 300 iterations. Starts included Dirichlet weights, sparse supports, uniform weights on random supports, and dense perturbations.

{\def\LTcaptype{none} % do not increment counter
\begin{longtable}[]{@{}rrrr@{}}
\toprule\noalign{}
Order & Starts & Convergence reports & Smallest \(\mathcal R\) \\
\midrule\noalign{}
\endhead
\bottomrule\noalign{}
\endlastfoot
8 & 40 & 40 & \(1-3.6\times10^{-15}\) \\
10 & 40 & 40 & \(1+7.1\times10^{-15}\) \\
12 & 40 & 40 & \(1+3.6\times10^{-15}\) \\
14 & 40 & 39 & \(1-7.1\times10^{-15}\) \\
16 & 40 & 38 & \(1\) \\
\end{longtable}
}

The values slightly below one occur near uniform weights and are at floating-point precision. They are not exact counterexamples.

Fixed-Gram searches kept the Fourier magnitudes and the real Nyquist eigenvalue fixed. They varied conjugate Fourier phases subject to nonnegative inverse-transform weights.

{\def\LTcaptype{none} % do not increment counter
\begin{longtable}[]{@{}
  >{\raggedleft\arraybackslash}p{(\linewidth - 8\tabcolsep) * \real{0.2000}}
  >{\raggedleft\arraybackslash}p{(\linewidth - 8\tabcolsep) * \real{0.2000}}
  >{\raggedleft\arraybackslash}p{(\linewidth - 8\tabcolsep) * \real{0.2000}}
  >{\raggedleft\arraybackslash}p{(\linewidth - 8\tabcolsep) * \real{0.2000}}
  >{\raggedleft\arraybackslash}p{(\linewidth - 8\tabcolsep) * \real{0.2000}}@{}}
\toprule\noalign{}
\begin{minipage}[b]{\linewidth}\raggedleft
Order
\end{minipage} & \begin{minipage}[b]{\linewidth}\raggedleft
Starts
\end{minipage} & \begin{minipage}[b]{\linewidth}\raggedleft
Convergence reports
\end{minipage} & \begin{minipage}[b]{\linewidth}\raggedleft
Smallest feasible ratio
\end{minipage} & \begin{minipage}[b]{\linewidth}\raggedleft
Smallest ratio among converged runs
\end{minipage} \\
\midrule\noalign{}
\endhead
\bottomrule\noalign{}
\endlastfoot
8 & 30 & 26 & 1.000000 & 1.000000 \\
10 & 30 & 27 & 1.120158 & 1.120158 \\
12 & 30 & 25 & 1.312790 & 1.312790 \\
14 & 30 & 22 & 1.096436 & 1.096436 \\
16 & 30 & 16 & 1.300263 & 2.533075 \\
\end{longtable}
}

Feasibility allowed minimum weight \(-10^{-7}\). Tiny negative weights were clipped for final ratio evaluation. Gram errors before clipping were below \(1.1\times10^{-14}\). The distinction between feasible outputs and converged outputs matters, especially at order sixteen. Initial and final weights and numerical residuals are retained for these runs.

\subsection{Configuration-model expectations}\label{app:finite-5}

For the uniform bipartite regular configuration multigraph of degree \(d\), the exact coefficient formulas in the main text evaluated

\[
\frac{c_n\,\mathbb E[\per(GG^{\T})]}{\mathbb E[\per(G)^2]}.
\]

Integer arithmetic checked all \(1\le n\le100\) for \(d=2,3,4\), and all \(1\le n\le60\) for \(d=5,10,20\). All 480 ratios were at most one. Equality occurred at order one. Selected numerical displays of these exact ratios are

{\def\LTcaptype{none} % do not increment counter
\begin{longtable}[]{@{}rrr@{}}
\toprule\noalign{}
\(d\) & \(n\) & Ratio \\
\midrule\noalign{}
\endhead
\bottomrule\noalign{}
\endlastfoot
2 & 100 & \(1.5610\times10^{-7}\) \\
3 & 100 & \(4.8180\times10^{-10}\) \\
4 & 100 & \(2.3821\times10^{-9}\) \\
5 & 60 & \(4.1872\times10^{-5}\) \\
10 & 60 & 0.0040689 \\
20 & 60 & 0.0576722 \\
\end{longtable}
}

Exhaustive stub-bijection enumeration independently checked the two expectation formulas at \((d,n)=(2,2),(3,2),(2,3)\). The respective counts were 24, 720, and 720. All sums agreed exactly. An expectation comparison does not give a pointwise inequality for every graph.

A separate stationary-point calculation used a Bethe approximation for the colored matching model at degrees two through twenty. It compared an estimated Gram logarithmic growth rate with \(1+2[(d-1)\log(d-1)-(d-2)\log d]\). It included a zero-message branch and six positive initial values for root finding. The largest recorded excess was approximately \(-0.048323685\). This is a numerical approximation. It does not establish the true limiting growth rate or a matrix inequality.

\subsection{Stable-polynomial and matroid deletion tests}\label{app:finite-6}

One suite had 52 examples at each degree three through six. At each degree, 40 were products of elementary-symmetric polynomials after a common doubly stochastic substitution. Twelve were determinant polynomials from balanced positive-semidefinite tuples. The initial coefficients and balance were exact; deletion balancing was numerical.

{\def\LTcaptype{none} % do not increment counter
\begin{longtable}[]{@{}
  >{\raggedleft\arraybackslash}p{(\linewidth - 6\tabcolsep) * \real{0.2500}}
  >{\raggedleft\arraybackslash}p{(\linewidth - 6\tabcolsep) * \real{0.2500}}
  >{\raggedleft\arraybackslash}p{(\linewidth - 6\tabcolsep) * \real{0.2500}}
  >{\raggedleft\arraybackslash}p{(\linewidth - 6\tabcolsep) * \real{0.2500}}@{}}
\toprule\noalign{}
\begin{minipage}[b]{\linewidth}\raggedleft
Degree
\end{minipage} & \begin{minipage}[b]{\linewidth}\raggedleft
Examples
\end{minipage} & \begin{minipage}[b]{\linewidth}\raggedleft
Smallest \(D_{\rm av}\)
\end{minipage} & \begin{minipage}[b]{\linewidth}\raggedleft
Smallest raw average ratio
\end{minipage} \\
\midrule\noalign{}
\endhead
\bottomrule\noalign{}
\endlastfoot
3 & 52 & 1.0000000 & 1.0000000 \\
4 & 52 & 1.0373282 & 1.0375230 \\
5 & 52 & 1.0242451 & 1.0242792 \\
6 & 52 & 1.0218714 & 1.0218926 \\
\end{longtable}
}

The largest balancing-gradient residual was below \(4.5\times10^{-8}\). All 208 recorded final ratios \(S\) were at least one.

A second suite had 90 examples per degree. Sixty used elementary-symmetric factors on different positive linear forms. Thirty used generic positive-semidefinite determinant pencils. Initial and deletion balancing used convex log-sum-exp minimization.

{\def\LTcaptype{none} % do not increment counter
\begin{longtable}[]{@{}
  >{\raggedleft\arraybackslash}p{(\linewidth - 6\tabcolsep) * \real{0.2500}}
  >{\raggedleft\arraybackslash}p{(\linewidth - 6\tabcolsep) * \real{0.2500}}
  >{\raggedleft\arraybackslash}p{(\linewidth - 6\tabcolsep) * \real{0.2500}}
  >{\raggedleft\arraybackslash}p{(\linewidth - 6\tabcolsep) * \real{0.2500}}@{}}
\toprule\noalign{}
\begin{minipage}[b]{\linewidth}\raggedleft
Degree
\end{minipage} & \begin{minipage}[b]{\linewidth}\raggedleft
Examples
\end{minipage} & \begin{minipage}[b]{\linewidth}\raggedleft
Smallest \(D_{\rm av}\)
\end{minipage} & \begin{minipage}[b]{\linewidth}\raggedleft
Smallest raw average ratio
\end{minipage} \\
\midrule\noalign{}
\endhead
\bottomrule\noalign{}
\endlastfoot
3 & 90 & 1.0198150 & 1.0198288 \\
4 & 90 & 1.0180926 & 1.0181045 \\
5 & 90 & 1.0204552 & 1.0204734 \\
6 & 90 & 1.0251072 & 1.0251346 \\
\end{longtable}
}

The largest residual was below \(2.4\times10^{-10}\). The ratio of the scaled deletion sum to the raw sum ranged from approximately 0.95592 to 0.999997. All 360 final ratios \(S\) passed. Determinant coefficients were obtained by integer finite differences with Bareiss elimination.

A third suite substituted positive linear forms into matroid basis polynomials, then balanced the result. Its 750 examples had the following minima.

{\def\LTcaptype{none} % do not increment counter
\begin{longtable}[]{@{}
  >{\raggedright\arraybackslash}p{(\linewidth - 8\tabcolsep) * \real{0.1579}}
  >{\raggedleft\arraybackslash}p{(\linewidth - 8\tabcolsep) * \real{0.2105}}
  >{\raggedleft\arraybackslash}p{(\linewidth - 8\tabcolsep) * \real{0.2105}}
  >{\raggedleft\arraybackslash}p{(\linewidth - 8\tabcolsep) * \real{0.2105}}
  >{\raggedleft\arraybackslash}p{(\linewidth - 8\tabcolsep) * \real{0.2105}}@{}}
\toprule\noalign{}
\begin{minipage}[b]{\linewidth}\raggedright
Source matroid
\end{minipage} & \begin{minipage}[b]{\linewidth}\raggedleft
Degree
\end{minipage} & \begin{minipage}[b]{\linewidth}\raggedleft
Examples
\end{minipage} & \begin{minipage}[b]{\linewidth}\raggedleft
Smallest \(D_{\rm av}\)
\end{minipage} & \begin{minipage}[b]{\linewidth}\raggedleft
Smallest \(S\)
\end{minipage} \\
\midrule\noalign{}
\endhead
\bottomrule\noalign{}
\endlastfoot
Fano & 3 & 150 & 1.0091356 & 1.0181519 \\
Dual Fano & 4 & 150 & 1.0213981 & 1.0648370 \\
Fano plus one coloop & 4 & 150 & 1.0167878 & 1.0513456 \\
Fano plus two coloops & 5 & 150 & 1.0190338 & 1.0782817 \\
Fano plus three coloops & 6 & 150 & 1.0149026 & 1.0762892 \\
\end{longtable}
}

These are Lorentzian inputs. The Fano source is not stable, although a substitution can produce a stable polynomial. Every raw average passed. The largest residual was below \(4.1\times10^{-9}\). Together, the three suites contain 1,318 saved examples. They do not prove a stable or Lorentzian deletion theorem.

Nine coefficient patterns obtained by grouping the seven Fano variables into three groups of at least two variables were also reported to pass numerical average deletion. No separate output record for those nine patterns is available. An additional 800 matrix-product tests, 200 at each order four through seven, reported no violation. Their individual inputs were not retained.

\subsection{Circulant deletion tests}\label{app:finite-7}

A coefficient-array implementation tested eight optimization starts at each order four through seven. All 32 runs reported convergence. The respective minima of \(D_{\rm av}\) were approximately \(1\), \(1+8.0\times10^{-11}\), \(1+1.7\times10^{-12}\), and \(1+5.3\times10^{-12}\). Residuals were below \(8.5\times10^{-13}\).

A permanent-gradient implementation agreed with the coefficient-array method within \(10^{-14}\) on 30 random cases in orders three through seven. It then tested 30, 30, 20, 9, and 6 starts at orders eight, ten, twelve, fourteen, and sixteen. All 95 runs reported convergence. No ratio below one beyond numerical precision was found. Capacity-gradient residuals were below \(2.8\times10^{-12}\). This numerical evidence is separate from the proved two-shift deletion theorem.

\subsection{Nested singleton deletion and concentrated elementary-symmetric inputs}\label{app:finite-8}

Nested tests started from a product of \(d+r\) positive linear forms in \(d+r\) variables. They extracted exponent one from \(r\) variables, then balanced the residual degree-\(d\) polynomial. Coefficient extraction used dynamic programming. Differentiation through capacity minimizers used covariance Hessians. Three small integer kernels agreed with naive coefficient expansion exactly. Directional gradient checks had errors below \(1.5\times10^{-9}\).

Average-deletion optimization saved 43 runs: four at \((d,r)=(3,3)\), seven at \((3,8)\), sixteen at \((4,6)\) across two parameterizations, eight at \((6,4)\), and eight at \((8,2)\). Forty-two reported convergence. No saved average or unrestricted maximum ratio was below one. Best-point residuals were below \(2.0\times10^{-12}\).

Prescribed-coordinate tests forced \(r+1\) factors to omit one variable, so its residual pure coefficient vanished. Twelve runs at \((3,5)\) had smallest prescribed ratio approximately 1.11622. One run each at \((4,0),(5,3),(6,2)\) gave approximately 0.95461005, 0.85057133, and 0.77161423. The average and unrestricted maximum ratios stayed above one. The exact order-four counterexample in the main text establishes failure of the prescribed rule. The degree-five and degree-six values remain numerical observations.

A compressed two-class calculation tested repeated-variable elementary-symmetric polynomials. Balance reduced to one noncentral-hypergeometric mean equation. Norms and deletions reduced to univariate coefficient products. Infinite multiplicity used a Poisson limit; multiplicity one used a balancing limit.

{\def\LTcaptype{none} % do not increment counter
\begin{longtable}[]{@{}rrr@{}}
\toprule\noalign{}
Degree & Cases & Smallest \(D_{\rm av}\) \\
\midrule\noalign{}
\endhead
\bottomrule\noalign{}
\endlastfoot
3 & 16 & 1.0033444425 \\
4 & 24 & 1.0025067642 \\
8 & 40 & 1.0012531916 \\
16 & 48 & 1.0006267742 \\
32 & 48 & 1.0003134584 \\
64 & 48 & 1.0001567502 \\
128 & 48 & 1.0000783808 \\
256 & 48 & 1.0000391918 \\
\end{longtable}
}

All 320 average, maximum, and final ratios passed numerically. The largest mean residual was approximately \(5.12\times10^{-12}\). Four small finite examples agreed with full multivariate expansion within \(6\times10^{-14}\). These tests do not prove an all-degree result.

\subsection{Direct Lorentzian coefficient optimization}\label{app:finite-9}

Cubic searches normalized the central coefficient and constrained balance, nonnegative coefficients, and derivative-Hessian signatures. Of 160 outputs, 134 passed the stated numerical feasibility screen: 63 of 80 deletion-objective outputs and 71 of 80 norm-objective outputs. The smallest recorded cubic deletion deficit \(G=3\sum_jH_j-4N\) among these outputs was approximately \(-3.76\times10^{-8}\). A numerical screen near a singular Hessian does not certify Lorentzianity or an exact negative deficit.

A boundary cubic search retained 62 outputs, of which 53 passed its screen. Its smallest \(G\) was approximately \(2/3\). A separate mixed-coefficient search retained 100 outputs, of which 91 passed. Its largest mixed sum was approximately \(10/3\). The exact polynomial in Appendix~\ref{app:additional-results} realizes these latter two values. The search does not prove the proposed universal boundary bound \(G\ge2b^2/3\).

At degree four, a 35-output boundary search had eight outputs that passed its numerical signature screen, but none reported optimizer convergence. Its smallest screened final ratio was approximately 1.003921125. A distinct eight-start direct constrained search had no feasible output. Those rejected outputs are not counterexamples.

A Schur-matrix barrier method then kept accepted iterates in the strict numerical Hessian region. It enforced balance through an affine nullspace. Four degree-four starts reached a smallest average ratio of approximately 1.003923425. Two degree-five starts reached approximately 1.00032812 and 1.00033093. The coefficient counts were 35 and 126; the Hessian counts were 10 and 35. Best-point balance residuals were below \(3\times10^{-13}\). The degree-four deletion-capacity residual reached \(2.05\times10^{-10}\) in one saved best point. These points have numerical feasibility checks only.

An additional cubic optimization was reported to have only nine feasible outputs among 120. Those nine gave no violation. Its underlying output file is not available, so its detailed constraints and values cannot be checked from a saved run.

Tests of fixed low-degree sum-of-squares forms used nonnegative combinations of Alexandrov--Fenchel minors, derivative-Hessian determinants, balance identities, and related reverse-Cauchy terms. Stable-only variants included Laguerre inequalities. The tested forms were reported infeasible. No exact dual certificate is supplied. This does not prove that other certificates are impossible. A numerical probe of the two-product determinantal relaxation also reported no failure, but no completed run record with counts and input values is available.

\subsection{M-convex support enumeration and attenuation}\label{app:finite-10}

For an M-convex support \(T\), let

\[
F_T(z)=\sum_{\alpha\in T}\frac{c_n}{\alpha!}
 e^{(\alpha-\mathbf1)\cdot z},\qquad \nabla F_T(z_*)=0.
\]

The target is \(F_T(2z_*)\le1\). A boundary support uses the corresponding balancing limit.

At degree three, exact subset enumeration and exchange checks found 57 labeled supports containing the center, with 17 coordinate-permutation orbits. All numerical target ratios passed. The largest balancing residual was approximately \(1.67\times10^{-12}\).

At degree four, the search covered all full-dimensional integral-polymatroid supports with the center in the interior. Of 5,184 candidate rank assignments, 615 were monotone and submodular. They gave 63 support orbits. Every support passed an exact exchange check. Numerical balancing gave no target violation. The largest proper-support value was approximately 0.9960941728. The largest residual was approximately \(9.63\times10^{-12}\). All 45,456 strict inclusion pairs among the 615 labeled supports were compared. No numerical inclusion-monotonicity violation was found.

A nonexhaustive degree-five beam search had width nine and a requested depth of 45. It completed depth levels one through thirty-three, visited 636 supports, and checked 1,537 rank-step inclusions. Rank and submodularity checks were exact. The largest proper-support value was approximately 0.9996800006. No target or inclusion failure was found. The largest state residual was approximately \(1.51\times10^{-8}\).

Fifty attenuation paths from asymmetric supports were sampled at seven interior parameters each. The largest interior value was approximately 0.6914288203. There were 252 samples above the endpoint chord. The largest interior balancing residual was approximately \(1.33\times10^{-12}\). The exact nonconvexity example in the main text establishes failure of continuous convexity independently of these numerical values.

The duplication comparison was also evaluated on the same 17, 63, and 636 supports. The complement-pair coefficient

\[
d_T=\sum_{\substack{\alpha\in T\\2\mathbf1-\alpha\in T}}
\frac1{\alpha!(2\mathbf1-\alpha)!}
\]

was exact. The denominator used numerical balancing. Every observed ratio passed, with minimum one at the full support. This does not prove a general duplication bound.

\subsection{Rank-two scalar and deletion tests}\label{app:finite-11}

A truncated-binomial scalar probe covered orders two through eighteen and every half-integer mean strictly between zero and the order. This gives 323 order-mean pairs. At each pair it enumerated admissible support intervals and solved each tilt equation numerically. It found no ratio below \(1-10^{-8}\). The smallest recorded ratio was 0.9999999999999953, at order eighteen and mean \(5/2\). The saved report contains the minimum and failures, not every interval-pair value. The associated broad ultra-log-concave assertion is not established by this probe; the exact obstruction in the main text must retain its stated scope.

For the shifted-binomial profile comparison, let \(B=\operatorname{Bin}(n,t/n)\), with \(0<t<n\). Let \(P,Q\) be shifted-binomial profile laws on \(\{0,\ldots,n\}\), both with mean \(t\). Write \(r=P/B\) and \(s=Q/B\) for their likelihood ratios. These laws satisfy \(0\le\operatorname{Var}(Q)\le t(1-t/n)\). The tested ratio was

\[
\frac{\langle r,s\rangle_B^2}{\langle r,T_{\kappa(s)}r\rangle_B},
\qquad
\kappa(s)=1-\frac{\operatorname{Var}(Q)}{t(1-t/n)}.
\]

Here \(T_\kappa\) is the binomial count-noise operator. It keeps each Bernoulli coordinate with probability \(\kappa\) and independently resamples it otherwise.

A numerical search of these profiles covered orders sixteen through twenty-two. It tested 175,574 ordered profile pairs and 5,723,130 initial trial means, followed by local one-dimensional refinements. It found no negative contraction deficit. The arithmetic was floating point. The trial schedule included points close to each endpoint and points inside each unit mean interval. Identical profiles and full-binomial profiles were omitted for separate mathematical reasons. This search does not extend the exact rank-two theorem beyond order fifteen.

Ten saved profile examples also violate a proposed endpoint-dominance shortcut. The exact interior-mean example in Appendix~\ref{app:additional-results} establishes that failure. A duplication-bridge diagnostic had a smallest recorded deficit of approximately \(-5.55\times10^{-16}\), at a full-binomial equality case. This value does not establish a strict negative deficit.

A rank-two deletion search saved an arbitrary-column failure at order three, with ratio approximately 0.96423013. The exact rational matrix in the main text proves that failure. Zero-containing-column searches saved extremizers at orders three through eight. Their smallest selected ratios were approximately

\[
1.12017221,\ 1.06856934,\ 1.05402263,\
1.03561156,\ 1.02709402,\ 1.02357208.
\]

A separate profile search took the maximum over the eligible zero-containing columns. Its six saved values, in the same orders, were approximately

\[
1.35637345,\ 1.12805402,\ 1.10598346,\
1.06057078,\ 1.04883878,\ 1.04642169.
\]

The saved extremizers and residuals permit inspection of those outputs. They are not a record of every sampled matrix and do not prove the zero-containing-column deletion assertion.

\subsection{Stochastic complements and minimum-norm deletion}\label{app:finite-12}

For a pivot \(a_{ij}<1\), define the stochastic complement

\[
B^{(ij)}_{k\ell}=a_{k\ell}+\frac{a_{kj}a_{i\ell}}{1-a_{ij}}
\quad(k\ne i,\ \ell\ne j),
\]

and let \(R(A)=\per(A)^2/\per(AA^{\T})\). The diagnostic pivot ratio is \(q_{ij}=g_nR(B^{(ij)})/R(A)\).

An initial matrix search retained extrema in orders three through seven. For each completed order, its program sampled 90 integer sums of permutation matrices. It tested duplication on the first 25 and tested all eligible pivots. The smallest duplication ratios ranged from 1.02143277 to 1.16495567. The largest observed minimum-pivot ratios were approximately 0.99540045, 0.89612551, 0.84375269, 0.87753059, and 0.89692640. Only the extrema were saved. The programmed order-eight run has no completed output.

A further pivot search retained order-three through order-seven summaries, with a schedule of 160 matrices per order. It examined minima and averages weighted uniformly, by entries, by complementary entries, by matching marginals, and by their complements. Some averages exceeded one. The largest uniform average among the retained summaries was approximately 1.10847118. The largest complement-of-matching average was approximately 1.39060221. The saved summaries contain their extremizing matrices. They do not establish the existence of a successful pivot for every matrix.

Direct constrained optimization maximized \(n^{-1}\sum_{ij}a_{ij}q_{ij}\). It saved 24 converged runs in each of orders three, four, and five, and 14 at order six. The largest scores were approximately one in the first three orders and 1.014299626 at order six. Constraint residuals were below \(3.6\times10^{-13}\). The exact stochastic-complement counterexamples in the main text establish the relevant failures without numerical optimization.

A separate search minimized the average of determinant-normalized minimum Fischer norms of deletions. It retained one extremizer in each order three through seven, from a schedule of 180 samples per order. The five average ratios were approximately 1.02219762, 1.02482586, 1.04882773, 1.03986009, and 1.03887917. Maximum balancing residuals were below \(7.1\times10^{-10}\). The exact order-700 example in the main text refutes this averaged minimum-norm comparison. These smaller numerical passes do not limit that counterexample.

\subsection{Finite arithmetic checks of analytic formulas}\label{app:finite-13}

Exact rational diagnostics checked 360 three-exceptional-row matrices in orders five through ten, and 51 at orders \(31,32,33,40,64\). The latter used a rank-three permanent expansion, checked against subset dynamic programming in orders five through seven. All tested margins passed. The all-order conclusions depend on the analytic proofs, not these samples.

Additional exact checks covered 405 integer-mean Poisson-binomial laws in orders two through ten, 28 two-column matrix identities, 184 four-block matrices in orders two through eight, and six row-pair affine identities. All passed. Separate local polynomial calculations checked the first and second variation formulas in degrees two through six for rank-one and full-rank projection directions. Those variation checks used numerical balancing and do not prove a neighborhood theorem.

For the finite rank-two certificate, a rational power-basis implementation agreed with an integer homogeneous implementation on all 454 profile intervals in orders two through six, comprising 20,452 coefficients. Direct permanent calculations tested 378 profile pairs in orders two through seven, with 1,390 identity comparisons, 21 singleton cases, and 251 uniform-profile instances. A further check supplied 209 direct comparisons. These tests support formula implementation. The exact whole-interval coefficient certificate and the analytic extremal reduction establish the finite theorem.

\subsection{Data limits}\label{app:finite-14}

The accompanying scientific data contain 69 computational JSON records and one JSON extraction from a numerical profile log. The extraction preserves every recorded numerical point and each completed-order count, including the printed decimal precision. These records retain input parameters or extremizing inputs, objective values, and residuals when those fields were available. Summary-only records cannot reconstruct every sampled point. The unavailable items identified above are the individual centered rank-two samples, the 800 additional matrix-product samples, the nine Fano-identification outputs, the 120-output cubic feasibility run, the sparse high-order matrix outputs, and completed output records for the two-product relaxation and fixed sum-of-squares probes.

The absence of those records does not alter an exact proof given in the main text. It limits what can be reproduced about the associated finite observation. No finite numerical pass, failed optimization, or unavailable diagnostic is used to assert the unrestricted Gyires inequality.
